\section{Amenazas a la validez}\label{sec:conclusiones:amenazas}

Esta sección agrupa por eje de validez las limitaciones que declaran los Capítulos~\ref{cap:catalogo} a~\ref{cap:comparacion}, cada una con su caso, su mitigación cuando la hay y lo que queda abierto.

\subsection{Validez de constructo}\label{sec:conclusiones:constructo}

La validez de constructo pregunta si lo medido es lo que se dice medir.
El sondeo de MailerSend midió el estado de la cuenta y no el plan: la cuenta sin aprobar admitió 10 peticiones por minuto frente a las 120 que el proveedor publica para una cuenta sana.
El sondeo de Amazon SES corrió en el entorno de pruebas (\textit{sandbox}) y el de Mailchimp Transactional en su cuenta de prueba.
Las tres Datasheets declaran el estado medido.
Medir en producción exige que el proveedor apruebe la cuenta, como en la salida del \textit{sandbox} de Amazon SES, o pagar un plan.

El \textit{rate} de Postmark y el de Mailgun salen de la cabecera de respuesta de una cuenta gratuita, sin provocar el código \texttt{429} con que la API rechaza el exceso, y la Datasheet los extiende a todos los planes.
En Postmark no lo provocaron 200 peticiones, y en Mailgun harían falta 80.000 en unos diez segundos.
Lo mitiga que las dos cifras vayan en ámbar, el color del semáforo para un valor interpretado y no literal.

SendGrid publica sus 10.000 peticiones por segundo en unas preguntas frecuentes y su envío no devuelve cabeceras de \textit{rate}, así que la cifra va en verde por ser literal, pero ningún sondeo la ha observado.
Queda abierta en las tres APIs una medición bajo carga, por ejemplo con el banco de pruebas de límites de uso de API-TaskForce, el grupo que publica la guía de ACM2YAML \citep{datasheetusagetesting}.

La curva de capacidad necesita la ventana de cada límite y lo que descuenta cada envío.
La ventana de la cuota, el periodo tras el que su contador vuelve a cero, está en rojo, el color del semáforo para un dato no publicado, en Mailgun, Brevo, Amazon SES y Mailchimp Transactional, y en ámbar en las otras cuatro.
Ninguna de las ocho declara la ventana del \textit{rate}, y los dos sondeos que la acotaron, en Brevo y Mailgun, no fijaron si es fija o deslizante.
Que cada destinatario descuente un correo de la cuota consta en Resend, SendGrid, Brevo y Amazon SES, y en Postmark, MailerSend, Mailgun y Mailchimp Transactional queda sin medir.
Esos huecos los marca el semáforo API a API, de modo que ninguno pasa por publicado sin serlo.

La cita literal de cada valor solo se contrasta de forma mecánica si se archivó la página, y de Postmark no se archivó ninguna de las que cita su Datasheet.
Lo mitigan la URL y la fecha de acceso de cada valor.

El ecosistema no tiene validador oficial de Datasheets, y la conformidad se comprueba con el validador de formato de este trabajo y con la ingestión en PRIME4API.
Los defectos de contrato, valores que ACM2YAML v0.3 admite y PRIME4API no sabe leer, los confirma un motor ajeno.
En cambio, las cuatro carencias que se promueven solo por recurrencia descansan en la lectura que este trabajo hace de la guía de ACM2YAML: los volúmenes contratados (P-01), el consumo de un lote de mensajes con varios destinatarios (P-02), el tope de una petición en otra unidad (P-03) y el estado de la suscripción (P-06).

P-06, el estado de la suscripción, muestra el riesgo de esa lectura.
La guía documenta alias de estado, bloques que dentro de un \textit{endpoint} dan a un estado de la cuenta sus propias cifras, con MailerSend como ejemplo, y la Datasheet de MailerSend añade uno para la cuenta sin aprobar.
En un alias, la guía solo documenta el \textit{rate} y la cuota, pero el código de referencia también lee un tope por petición.
El tope de dos destinatarios por petición medido en esa cuenta cabe, por tanto, en el esquema, y lo que falta es documentarlo en la guía.
P-06 sigue siendo una extensión porque el esquema no permite declarar el disparador del estado ni declarar el estado una sola vez para toda la Datasheet.

A cada pregunta se le asigna el instrumento que la responde, sea la Datasheet, PRIME4API o HARVEY4API, sabiendo ya lo que ofrece cada uno, y eso puede inflar la proporción de preguntas respondibles.
Por eso el 78~\% de preguntas respondibles, 25 de las 32 del catálogo de preguntas, no se compara de forma directa con el 54~\% que Cavero et al.\ obtienen sobre 56 plantillas con herramientas que ya existían \citep{cavero2026pricingintelligence} (Sección~\ref{sec:comparacion:catalogo}).
Lo mitiga que las preguntas sigan la gramática de plantillas de Cavero et al., de modo que pueden incorporarse a su banco de preguntas y medirse con sus herramientas.

El catálogo de preguntas se cerró en 32 sin un criterio de saturación fijado de antemano, aunque al pasar de tres APIs a ocho no ganó ninguna.

\subsection{Validez interna}\label{sec:conclusiones:interna}

La validez interna pregunta si los resultados dependen de las Datasheets o del montaje con que se obtienen.
El script que verifica las Datasheets contra PRIME4API clona el motor sin fijar la revisión, y un cambio del motor puede mover los errores de la Tabla~\ref{tab:verificacion} sin tocar la Datasheet.
Seis de las ocho ejecuciones registran la revisión \texttt{8e6011c} del repositorio de PRIME4API y pueden repetirse contra ella, y las de Brevo y Amazon SES no la registran.

El método de sondeo tiene tres puntos débiles.
En MailerSend, el \textit{endpoint} de cuota servía el saldo del día anterior once horas después de la reposición (Sección~\ref{sec:limites:catalogo}), así que el saldo que registra el sondeo al terminar puede ser el de la víspera.
En Brevo, el sondeo lee la cifra en el error con que el proveedor rechaza un destinatario imposible de entregar, pero con la cabecera de modo \textit{sandbox}, que solo valida el formato, el envío devolvió \texttt{201}.
Lo cubren un sondeo previo que comprueba que el proveedor rechaza la petición y un tope de envíos que cuenta como real cada petición aceptada.
En Mailchimp Transactional, desde once destinatarios o con un adjunto de 5~MB, el envío queda en cola, y por encima de ese umbral el sondeo no mide.

PRIME4API corrió en local con la misma imagen de contenedor que despliega SPHERE, y no en el despliegue público que cita Gavira Sánchez \citep{gavira2026tfm}, así que sus respuestas no se han contrastado contra ese servicio.
Las diez respuestas de Mailchimp Transactional que se calcularon a mano por el HTTP 500 del motor (Sección~\ref{sec:conclusiones:evaluacion}) no tienen confirmación de PRIME4API.

Los valores modelados dependen también de qué fuente se toma cuando dos discrepan.
La regla que da preferencia a la página de precios vale para las 37 divergencias, de las 82 registradas, en que esa página es uno de los lados.
Aun en ellas deja sin resolver las siete en que la página se contradice a sí misma, cuyo valor queda en ámbar.
Las 45 restantes enfrentan fuentes distintas de la página de precios: en la mayoría, dos fuentes del mismo proveedor, como su referencia de la API y su centro de ayuda, y en las demás, una de ellas frente a un sondeo o a un tercero.
Las ordenan otras reglas de modelado, como la que da preferencia al valor medido.
Dieciséis de esas 45 siguen abiertas, con las dos fuentes en el registro de procedencia.

Las Datasheets y el recuento de recurrencia los construye un agente de inteligencia artificial, que podría dar por leído un dato que la fuente no contiene.
Lo mitiga que el agente que construye no sea el que verifica, y queda abierto medir cuántos valores corrigió esa revisión (Sección~\ref{sec:conclusiones:lecciones}).

\subsection{Validez externa}\label{sec:conclusiones:externa}

La validez externa pregunta hasta dónde valen los resultados fuera del catálogo y de su ventana de toma de datos.
La respuesta negativa a la RQ2, si ACM2YAML v0.3 basta para representar las APIs de correo, no depende de la muestra, porque basta un caso que v0.3 no represente, como los seis volúmenes contratados de Postmark.

Sí depende de la muestra qué propuestas se promueven al esquema.
El tope de una petición medido en una unidad distinta de la suya, como los 50~MB de carga de \texttt{/email/bulk} en Postmark (P-03), aparece en las ocho APIs.
En cambio, tres propuestas descansan en una sola API: el tope del sobrecoste de SendGrid (P-07), el cupo de destinatarios de la cuenta sin aprobar de MailerSend (P-08) y la cola de Brevo por encima de la cuota (P-09).
Se promueven porque PRIME4API da en ellas un resultado que el proveedor contradice.
Al pasar de cuatro APIs a ocho ninguna propuesta perdió su apoyo.
Queda abierto repetir el inventario en otro dominio, comprobando cada propuesta con el caso que la motiva.

La marca de valor no publicado (P-04) y el plan con precio no público (P-05) se han medido contra un único motor, PRIME4API.
Claude Opus 5.5, con el que se contrastaron las preguntas de HARVEY4API, lee el YAML como texto y no cuenta como segundo motor.

En la relación entre límites de uso y \textit{rates} (RQ3), que la cuota sea la restricción activa se explica por las unidades, y por eso cabe esperarlo fuera del catálogo: frente a una cuota mensual o diaria, un \textit{rate} por segundo o por minuto, sostenido durante la ventana de la cuota, suma mucho más volumen.
Lo que sí depende de la muestra son las excepciones, como el \textit{rate} reducido de la cuenta sin aprobar de MailerSend (Sección~\ref{sec:conclusiones:constructo}) o la discrepancia de Mailgun.
Mailgun publica 1.500 peticiones por segundo por IP de cliente, y la cabecera del envío devuelve un límite de 80.000 peticiones en una ventana medida de unos diez segundos.
Que el \textit{rate} del envío individual no cambie con el plan vale en Resend, SendGrid y MailerSend, y Brevo lo contradice.

Las cuatro formas de cobro que el catálogo cubre (volumen mensual, pago por uso, producto añadido y límite diario) salen del primer criterio de selección, cubrir formas de cobro distintas (Sección~\ref{sec:propuesta:criterios}), y no de un censo del sector.
Entre las APIs de reserva de la selección, Zoho ZeptoMail cobra con créditos prepagados, una variante que el catálogo no incluye.

El catálogo modela solo el producto transaccional, el de los mensajes individuales que suele desencadenar el propio usuario, y sus conclusiones no se trasladan al de \textit{marketing}, que Resend factura por contactos y no por correos.

La oferta cambia con el tiempo: el repositorio de ACM2YAML conserva una Datasheet de MailerSend con precios que difieren de los publicados cinco meses después \citep{apisconsumptionmodel}.
Cada valor lleva su fecha de acceso, y falta un procedimiento que diga cuándo reverificar.

\subsection{Validez de conclusión}\label{sec:conclusiones:conclusion}

La validez de conclusión pregunta si los recuentos que sostienen las conclusiones están bien hechos.
La Tabla~\ref{tab:recurrencia} deja 31 combinaciones de elemento y API sin verificar y una en suspenso, las dos cuotas de Mailgun que ninguna página dice si se aplican a la vez.
Una aparición sin verificar no se cuenta, así que verificarla solo puede añadir APIs y cada recuento es un mínimo.
Bastaría confirmar una aparición más en cada uno para que pasaran a recurrentes tres elementos no promovidos: la región como unidad de gobierno, el tope fraccionario y los dos contadores sobre una misma ventana.
Ninguna promoción se retiraría\footnote{El detalle por fila está en \texttt{03\_datos/\allowbreak RECURRENCIA\_\allowbreak 2026-09-23.md}, en el repositorio del trabajo.}.

En Postmark, una respuesta «no publicado» cuenta como respondida en la pregunta del \textit{rate} máximo y no en la del límite diario ni en la de si el \textit{rate} sube con el plan.
Aplicar el mismo criterio a las tres movería una o dos de sus 18 respuestas.
La cobertura cuenta igual el ámbar que el verde, y el ámbar lo llevan 13 de las 18 respuestas de Postmark y 12 de las 20 de Mailgun, frente a 1 de 22 en Resend.

El mes se toma de 30 días.
Con 28, el menor cociente entre \textit{rate} y cuota de una cuenta sana, el de MailerSend Professional 2M, un volumen que la Datasheet no recoge, baja de 2,59 a 2,42.
Ninguna restricción activa cambia con la longitud del mes.
